\begin{table}[H]\centering\scriptsize\caption{Caracterización preliminar de las fases CIP.}
\begin{tabularx}{\textwidth}{L{1.9cm}L{2.3cm}L{1.4cm}L{2.0cm}YL{1.4cm}L{1.6cm}}\toprule
\textbf{Fase}&\textbf{Fluido}&\textbf{Temp.}&\textbf{Tiempo}&\textbf{Bomba/válvulas}&\textbf{Transición}&\textbf{Destino}\\\midrule
Preenjuague&Agua potable, 15--20 L, máx. 20 L (paso único)&50\,$^{\circ}$C (preliminar)&2--3 min (preliminar; se define según el tiempo de paso del líquido total por las tuberías)&Bomba de impulsión encendida; bomba dosificadora apagada; válvula de salida a drenaje abierta&Sin cambios&Drenaje\\
Lavado&Agua potable, 15--20 L, máx. 20 L, con detergente diluido al 1,5\,\% v/v, aprox. 225--300 mL (paso único)&40\,$^{\circ}$C (preliminar)&5--10 min (preliminar; se define según el tiempo de paso del líquido total por las tuberías)&Bomba de impulsión encendida; bomba dosificadora activa solo durante la dosificación; válvulas dirigen el fluido en un solo paso hacia el equipo de prueba y de ahí al drenaje&Sin cambios&Drenaje\\
Enjuague final&Agua potable, 15--20 L, máx. 20 L (paso único)&Ambiente&2--3 min (preliminar; se define según el tiempo de paso del líquido total por las tuberías)&Bomba de impulsión encendida; bomba dosificadora apagada; válvula de salida a drenaje abierta&Sin cambios&Drenaje\\\bottomrule
\end{tabularx}\end{table}
